% =====================================================================
%  Aufgabenstellung (konsolidiert) — DEUTSCH (treue Uebersetzung)
%  Konsolidiert aus den Projektterminen docs/termine/ (Termin 1-9) UND der
%  gesamten Architektur-Doku (docs/architektur/14 Organ-Metapher, Bausteine-
%  Matrix N, Kapitel 01/04). Original (Quelle): en.tex.
% =====================================================================
{\large\textbf{Aufgabenstellung}}\\[0.4em]
{\itshape Aktive cache-bewusste Hardware-Adaption:\\ Eine Cache-Engine für Trie-basierte Indexstrukturen}

\bigskip
\minisec{Kontext und Motivation}
Trie-basierte Indexstrukturen sind ein Kernbaustein moderner
In-Memory-Datenbanksysteme, und die Professur für Datenbanken adressiert
ausdrücklich die hardware-nahe Optimierung solcher Systeme. Seit dem Adaptive
Radix Tree wurde eine Vielzahl cache-bewusster Techniken---Knoten-Layout,
Präfixkompression, Prefetching, Succinct-Repräsentationen und Allokation---vorgeschlagen,
doch nutzt jede Veröffentlichung eigene Datenstrukturen, Lasten und
Messmethodik. Es fehlt ein gemeinsames Rahmenwerk, das diese Entwürfe in
austauschbare Bausteine zerlegt und auf einer einzigen, fairen Basis misst.

\minisec{Zielsetzung}
Die Arbeit soll ein Achsen-Bibliotheks-Framework (eine \emph{Cache-Engine})
entwerfen und prototypisch umsetzen, das cache-bewusste Trie-Indizierung in
orthogonale, unabhängig austauschbare Bausteinachsen zerlegt, Algorithmen des
Stands der Technik als explizite Achsen-Konfigurationen rekonstruiert und
sowohl einzelne Achsen als auch ganze Algorithmen auf gemeinsamer, CPU-only
Basis misst. Als Prüfling für die Zugehörigkeit zum Stand der Technik trägt die
Arbeit \mbox{PRT-ART} bei (einen Probst-Redirect-Tree-/ART-Hybriden).

\minisec{Methodischer Ansatz: Algorithmen als ``Tiere'', Achsen als ``Organe''}
Das Framework beruht auf einer anatomischen Metapher. Ein Suchalgorithmus wird
als \emph{Tier} (Lebewesen) aufgefasst; die Teilaufgaben, die \emph{jeder}
Algorithmus erfüllen muss, sind seine \emph{Organe}. Jedes Organ ist eine
orthogonale \emph{Achse}: So wie jeder Wiederkäuer einen Magen besitzt, ihn aber
Kuh, Reh und Schaf unterschiedlich ausprägen, besitzt jeder Suchalgorithmus ein
Knoten-Layout-Organ, ein Traversal-Organ, ein Allokations-Organ usw., jeweils in
anderer Ausprägung. Ein konkreter Algorithmus ist dann \emph{eine Komposition}
--- ein einzelner Punkt im kartesischen Produkt aller Achsen; ein Entwurf des
Stands der Technik (ART, HOT, Masstree, \ldots) ist genau ein solcher Punkt. Die
Permutation der Achsen ist ein \emph{genetisches Experiment}: Der Austausch von
Organen zwischen Tieren erzeugt neue Komposit-Algorithmen mit bisher unbekannten
Eigenschaften. Das Forschungsziel ist es, die gemeinsame \emph{Anatomie} eines
Suchalgorithmus zu finden und im resultierenden Permutationsraum das schnellste
``Tier'' --- die optimale cache-line-bewusste Komposition --- für ein gegebenes
datentyp-vorzügliches Lastmuster zu bestimmen.

\minisec{Bausteinachsen und Cache-Line-Relevanz}
Eine Teilaufgabe der Arbeit ist zu ermitteln, welche Achsen (Organe) jeder
Suchalgorithmus besitzt. Die konsolidierte Matrix umfasst \textbf{vierzehn
Hauptachsen plus etwa dreißig Sub-Achsen}: (1)~Page-Type, (2)~Node-Type,
(3)~Traversal --- aufgeteilt in 3.A~Algorithmus-Seite, 3.B~Cache-Memory-Seite
und 3.M~Mapping ---, (4)~ValueHandle, (5)~Memory-Layout, (6)~Allocator
(Allokation, Reclamation, NUMA, Huge-Page, Free-List), (7)~Prefetch,
(8)~Concurrency (Pattern, Locking-Mode), (9)~ISA, (10)~Measurement,
(11)~Telemetry, (12)~Hardware-Strategie, (13)~Scheduling-Strategie und
(14)~Identity. Die \textbf{cache-line-bewusst relevanten} Organe sind
insbesondere: Node-Type und Memory-Layout (wie Knoten 64-Byte-Cache-Lines
belegen), die \emph{Cache-Memory-Traversal} (3.B) zusammen mit ihrem
\emph{Mapping} (3.M, das die nicht-cache-kohärente Algorithmus-Sicht des
Entwicklers in die tatsächliche Cache-Mechanik auf der Speicherseite übersetzt),
der Allocator (NUMA-Affinität und Huge-Pages bestimmen Cache-/TLB-Lokalität),
Prefetch, ISA und die Hardware-Strategie.

\minisec{Warum die vollständige Permutation erforderlich ist}
Cache-line-bewusstes Verhalten entsteht nicht aus einem einzelnen Organ, sondern
aus dem \emph{Zusammenspiel} vieler --- Knoten-Layout $\times$ Memory-Layout
$\times$ Allocator $\times$ Prefetch $\times$ Traversal-Mapping $\times$ ISA. Ein
Organ lässt sich nicht isoliert cache-line-bewusst optimieren; die
Wechselwirkungen zwischen den Achsen sind entscheidend. Die Arbeit misst den
Konfigurationsraum daher \emph{vollständig und permutativ}, ohne Reduktion
(``Messung immer so vollständig wie möglich''); das kartesische Produkt spannt
mehr als~$10^{11}$ Permutationen auf, die der Builder als einzelne, ABI-stabile
Mess-Binaries erzeugt.

\minisec{Teilaufgaben}
\begin{enumerate}
  \item Systematisierung der Organe cache-bewusster Trie-Indizes zu einer
        kanonischen Permutationsmatrix orthogonaler Achsen und Identifikation der
        cache-line-bewusst relevanten Achsen.
  \item Rekonstruktion von Entwürfen des Stands der Technik (ART, HOT, Masstree,
        CoCo-trie, START, B\textsuperscript{2}-Tree, Wormhole, SuRF sowie
        weitere Tier-2/3-Entwürfe) als explizite Achsen-Kompositionen
        (Referenz-``Tiere'').
  \item Entwurf und Implementierung von PRT-ART als Prüfling mit eigenen Organen
        (Redirect-Knoten, adaptive lokale Suchseiten, externe Value-Arenen,
        cache-line-bewusstes Layout).
  \item Aufbau einer reproduzierbaren Mess-Pipeline
        (Binary~$\to$~CSV~$\to$~LaTeX~$\to$~Diagramm) mit profilbewusstem
        YCSB-Workload-Routing und Hardware-Counter-Erfassung.
  \item Vollständige Durchführung der drei verpflichtenden Messreihen:
        (A)~Prüfling vs.\ Stand der Technik, (B)~Cache-Engine-Permutationen
        (Allokator~$\times$~Layout~$\times$~Prefetch~$\times$~\ldots),
        (C)~Merge/Regression alter vs.\ neuer Varianten.
  \item Identifikation der schnellsten cache-line-bewussten Kompositionen je
        datentyp-vorzüglichem Lastmuster aus dem vollständigen Permutationsraum
        (Forschungsziel).
\end{enumerate}

\minisec{Methodik und Evaluation}
Die Evaluation ist CPU-only und vergleicht die Kompositionen gegen etablierte
Baselines unter variierten Datenvolumina, Keylängen, Präfixüberlappungen,
lokaler Dichte, Skew, Treffer-/Fehlerquoten und Value-Größen (den
datentyp-vorzüglichen Lastmustern). Gemessen werden Exact-/Prefix-Lookup,
Prefix-Enumeration, Insert/Update/Delete, p50/p95/p99-Latenz, Durchsatz, Bytes
pro Key sowie Hardwarezähler (cycles, instructions, L1/L2/LLC- und dTLB-Misses,
Branch-Mispredictions). Zunächst wird Single-Thread-Verhalten untersucht, danach
folgen leseparallele Multi-Thread-Szenarien.

\minisec{Umfangsabgrenzung}
Der Kernbeitrag ist CPU-only und single-threaded mit Lese-Skalierung. GPU liegt
außerhalb des Umfangs; Schreibparallelität, SIMD und Engine-/Optimierer-Integration
sind optionale Erweiterungen.
